%======================================================================
\section{Introduction}
\label{sec:intro}
%======================================================================
Spark jobs are trusted because they are boring. A transformation written against the DataFrame API does the same thing on every run, its logical plan can be inspected before execution, and a failed stage can be retried without changing the result~\cite{zaharia2012rdd,zaharia2016spark,armbrust2015sparksql}. Data cleaning has always sat awkwardly inside this model. Removing exact duplicates or trimming whitespace is mechanical, but deciding that \texttt{ENGINEERING} and \texttt{engineering} denote the same department, that \texttt{15/01/2022} should be rewritten as \texttt{2022-01-15}, or that a missing salary should be imputed at all, requires judgement about what the data means. In practice that judgement is either hard-coded as rules or left to an analyst~\cite{rahm2000cleaning,chu2016cleaning,ilyas2019cleaning}.

Large language models (LLMs) can supply some of this judgement from a short description of a column~\cite{narayan2022foundation,zhang2023preprocessors}. The difficulty is what happens when they are wrong. A model that proposes a plausible but unsupported repair, refers to a column that does not exist, or misreads a statistic~\cite{ji2023hallucination,huang2025hallucination} produces exactly the kind of error that a pipeline cannot detect after the fact: the job succeeds, the output has the right schema, and the damage is spread across millions of rows. Integration patterns that let a model write and run code, such as text-to-SQL and tool-calling agents~\cite{rajkumar2022texttosql,li2023bird,yao2023react,schick2023toolformer}, make this failure mode easier to reach.

This paper evaluates one answer to the question in its title. CONVERGE is a pipeline in which the model is allowed to \emph{recommend} cleaning operations but never to execute them. The model receives a bounded summary of the data computed by Spark and returns typed recommendations drawn from a closed catalog of eight operations. A deterministic validator, which has no dependency on the model, re-reads the same Spark evidence and decides which recommendations may run. Spark then executes the approved plan unchanged and refuses to publish partial output.

We compare this design with two alternatives that share every other component: a Spark-only baseline that applies the two repairs the evidence fully determines, and an \textsc{AI\_Only} pipeline that executes the same model's recommendations after a structural filter. Because the profiler, prompt, executor and evaluator are identical across the three arms, differences in outcome are attributable to the decision layer.

The study makes five contributions.
\begin{enumerate}[leftmargin=*,nosep]
  \item An architecture for model-assisted cleaning in Spark with five stated safety properties, implemented as a ten-state auditable pipeline (\S\ref{sec:problem}, \S\ref{sec:arch}).
  \item A controlled-corruption benchmark with per-cell ground truth and metrics that separate repair quality, collateral damage and run availability from the data-quality signals a pipeline normally observes (\S\ref{sec:method}).
  \item An end-to-end evaluation with a locally served 3B-parameter model across six scenarios and three corruption seeds, with bootstrap intervals and paired permutation tests (\S\ref{sec:res-e1}, \S\ref{sec:res-seeds}).
  \item A record-and-replay methodology that re-runs the model's actual outputs through fifteen decision pipelines, giving paired, causal estimates of what each validator rule contributes and evaluating three defect fixes and two new semantic guards on held-out seeds (\S\ref{sec:res-replay}).
  \item A twelve-class hallucination taxonomy with an injection study, a token-level cost analysis, and a scalability study to $10^6$ rows, which together locate precisely what validation can and cannot prevent (\S\ref{sec:res-inject}, \S\ref{sec:res-scale}).
\end{enumerate}

The headline result is mixed in an informative way. The model was necessary for one class of repair (date formats, recall $0\to1$) and harmful for another: on the missing-value and categorical scenarios it invented problems and rewrote correct values, and the comparator that a pipeline would normally use to judge cleaning reported those runs as complete successes. The shipped validator removed most of that damage, but the replay study shows that part of its apparent success came from a defect, and that two small deterministic guards are needed before the design behaves as intended on unseen data.

%======================================================================
\section{Background and Related Work}
\label{sec:related}
%======================================================================
\subsection{Data Cleaning at Scale}
Classical cleaning separates error detection from repair~\cite{rahm2000cleaning,chu2016cleaning,ilyas2019cleaning,hellerstein2008quantitative}. Constraint-based approaches such as conditional functional dependencies~\cite{fan2008cfd} give repairs a formal meaning but need curated rules. HoloClean~\cite{rekatsinas2017holoclean} combines constraints, statistics and external signals in a probabilistic model; ActiveClean~\cite{krishnan2016activeclean} and Raha~\cite{mahdavi2019raha} reduce the human effort needed to configure detection. Abedjan et al.~\cite{abedjan2016detecting} found that no single detector dominates across real datasets. Closer to pipeline engineering, Deequ~\cite{schelter2018deequ} expresses data-quality checks as Spark computations, and TFX data validation~\cite{breck2019validation} checks training data against a schema before it reaches a model. CONVERGE borrows the second idea, deterministic checks computed by the engine itself, but applies it to \emph{proposed actions} rather than to the data alone.

\subsection{Language Models for Data Preparation}
Narayan et al.~\cite{narayan2022foundation} showed that foundation models can match specialised systems on entity matching, imputation and error detection with few examples, and Zhang et al.~\cite{zhang2023preprocessors} reported similar results for a range of preprocessing tasks. Earlier, Wrangler~\cite{kandel2011wrangler} established the value of \emph{suggesting} transformations to a user rather than applying them. These studies measure whether the model's answer is right. They do not ask what a pipeline should do when it is wrong, which is the question for any unattended deployment. Text-to-SQL benchmarks~\cite{rajkumar2022texttosql,li2023bird} and agent frameworks~\cite{yao2023react,schick2023toolformer,shinn2023reflexion} go further by executing model output; our \textsc{AI\_Only} arm reproduces that pattern under a restricted action space.

\subsection{Hallucination and Guarding Model Output}
Surveys of hallucination~\cite{ji2023hallucination,huang2025hallucination} distinguish errors that contradict the input from errors that contradict the world. Retrieval augmentation~\cite{lewis2020rag} targets the second; constrained decoding~\cite{willard2023outlines} guarantees that output is syntactically valid but says nothing about whether it is true. Indirect prompt injection~\cite{greshake2023injection} shows that anything the model reads can carry instructions, which matters here because evidence contains raw cell values. From a security standpoint, Saltzer and Schroeder's principles of least privilege and complete mediation~\cite{saltzer1975protection} argue that an untrusted component should hold no execution authority and that every request it makes should be checked. CONVERGE applies these principles literally: the model can only propose, and every proposal passes through a validator that does not depend on the model.

\subsection{Positioning}
Existing work either cleans deterministically without semantic judgement or uses models for judgement without an execution boundary. We are not aware of a study that measures, against per-cell ground truth, how an evidence-based validation boundary changes the outcome of model-driven cleaning, or one that decomposes that change rule by rule on the model's own outputs.

%======================================================================
\section{Problem Formulation and Threat Model}
\label{sec:problem}
%======================================================================
Let $D$ be a table with columns $C=\{c_1,\dots,c_m\}$ and $n$ rows. The Spark profiler computes $\pi(D)=(P,E)$, where $P$ holds per-column statistics (physical and semantic type, null and distinct counts, quantiles, date-pattern histograms) and $E=\{e_1,\dots,e_k\}$ is a set of evidence items. Each $e=(\tau,C_e,\sigma)$ names an issue type $\tau$ (missing values, duplicates, categorical inconsistency, whitespace, date-format inconsistency, outlier candidate), a column set $C_e\subseteq C$, and bounded supporting data $\sigma$ such as a variant histogram. The sizes of $E$ and $\sigma$ are capped by configuration, not by $n$.

The operation catalog $\mathcal{O}$ assigns each permitted operation $o$ a safety tier $t(o)\in\{\textsc{safe},\textsc{review},\textsc{finding}\}$, compatible semantic types $S(o)$, a parameter schema $\Theta(o)$ and required evidence types $R(o)$. A separate set $\mathcal{F}$ of forbidden names (e.g., \texttt{DELETE\_COLUMN}, \texttt{ARBITRARY\_CODE}, \texttt{JOIN}) is recognised and always refused.

The model $M$ maps $(P,E,\mathcal{O})$ to recommendations $r=(o,C_r,\theta,\rho,\mathrm{refs})$: an operation, target columns, parameters, a free-text reason, and evidence references with \emph{claimed values}. We treat every field of $r$ as untrusted.

\noindent\textbf{Checkable hallucinations.} A recommendation is hallucinated if it violates any of the following: (H-op) $o\in\mathcal{O}\setminus\mathcal{F}$; (H-col) $C_r\subseteq C$; (H-type) every $c\in C_r$ has $\mathrm{sem}(c)\in S(o)$; (H-param) $\theta$ conforms to $\Theta(o)$ with no extra keys; (H-evid) some $e\in E$ has $\tau_e\in R(o)$ and covers $C_r$ (with equality for row-level operations); (H-claim) each claimed value equals the corresponding entry of $\sigma$; (H-value) each raw value named in $\theta$ occurs in $\sigma$. These predicates can be evaluated from $(P,E)$ alone. A recommendation can also be \emph{semantically wrong} while satisfying all of them, for instance by mapping a genuine variant of ``Marketing'' to ``Sales''. Part of our evaluation asks how much of this second class can be recovered by additional predicates that are still computable from evidence (\S\ref{sec:res-replay}).

\noindent\textbf{Safety properties.} The design aims to guarantee:
\begin{enumerate}[label=\textbf{P\arabic*},leftmargin=*,nosep]
  \item \emph{Bounded exposure}: $M$ observes only $(P,E,\mathcal{O})$; never rows, file paths or engine handles.
  \item \emph{Closed action space}: outputs are typed data, and no field is evaluated as code.
  \item \emph{Complete mediation}: no $r$ is executed unless the validator returns an executable status.
  \item \emph{Monotone approval}: an approval policy can narrow the validator's verdicts but not widen them.
  \item \emph{Immutable, fail-fast execution}: the executor applies a frozen plan in order, re-checks parameters, and writes nothing if any operation fails.
\end{enumerate}
The threat model is a model that is mistaken rather than adversarial, although P1--P3 also bound the effect of injected instructions carried in cell values.
